% ----------------------------------------------------------------------------
%                                   RESUMO
% ----------------------------------------------------------------------------

\begin{resumo}[RESUMO]

    \begin{SingleSpacing}
    
        Este trabalho apresenta o desenvolvimento de um sistema web voltado ao auxílio no manejo comportamental de crianças diagnosticadas com Transtorno de Déficit de Atenção com Hiperatividade (TDAH) e Transtorno Opositor Desafiador (TOD). O objetivo central é promover a comunicação integrada entre pais e professores, de modo a alinhar estratégias de acompanhamento e reduzir inconsistências no processo educativo. A pesquisa parte da constatação de que a falta de comunicação eficaz entre família e escola compromete o desempenho acadêmico e social dessas crianças, além de gerar conflitos e dificultar a aplicação de intervenções consistentes. Quanto à metodologia científica, o estudo classifica-se como uma pesquisa tecnológica e aplicada. Para estruturar a solução, adota-se a prototipação como método de desenvolvimento, precedida pelo uso de instrumentos de pesquisa qualitativos — como a aplicação de questionários, entrevistas semiestruturadas e observação participante com a comunidade escolar — essenciais para o levantamento preciso de requisitos. A partir disso, o projeto busca oferecer funcionalidades para o registro ágil de comportamentos e o acompanhamento de estratégias pedagógicas proativas. A avaliação da ferramenta será conduzida posteriormente por meio de testes de usabilidade com usuários finais, garantindo a eficiência da interface. Os resultados esperados incluem a melhoria da comunicação entre os atores envolvidos, a redução da sobrecarga docente e o fortalecimento dos vínculos sociais, reforçando a importância da tecnologia voltada à inclusão e ao apoio no desenvolvimento infantil.
        
        \item\textbf{Palavras-chave}: Sistema web. Manejo comportamental. TDAH. TOD. Prototipação.
    
    \end{SingleSpacing}
    
\end{resumo}
